\documentclass[12pt]{article}
\usepackage[margin=1in]{geometry}
\usepackage{times}
\usepackage[T1]{fontenc}
\usepackage[utf8]{inputenc}
\usepackage{hyperref}
\usepackage{tcolorbox}
\usepackage{enumitem}
\usepackage{titlesec}
\usepackage{parskip}
\usepackage{setspace}

\hypersetup{colorlinks=true, linkcolor=blue, urlcolor=blue, citecolor=blue}
\tcbuselibrary{skins}

\newtcolorbox{callout}[1]{
  colback=gray!8, colframe=gray!40, boxrule=0.4pt,
  arc=1mm, left=3mm, right=3mm, top=2mm, bottom=2mm,
  fonttitle=\bfseries, title=#1
}

\titleformat{\section}{\large\bfseries}{\thesection}{0.5em}{}
\titleformat{\subsection}{\normalsize\bfseries}{\thesubsection}{0.5em}{}

\title{\Large\bfseries The Corporate Stage: Algorithmic Surveillance, Hybrid\\
Visibility, and the Paradox of Spectatorship in the Modern\\
Workplace}

\author{
Vivien Jiaqian Zhu\\
\small Department of East Asian Languages and Cultures\\
\small University of California, Berkeley\\
\small Visiting Scholar, Stanford University\\
\small Research Fellow, University of Cambridge\\
\small \texttt{email@example.com}
}
\date{}

\begin{document}
\maketitle

\begin{abstract}
\noindent The shift to hybrid and remote work has coincided with the normalization of continuous algorithmic monitoring --- keystroke logging, activity scoring, webcam-based presence detection --- producing a workplace in which supervision and self-presentation now operate under fundamentally altered conditions. This chapter develops a theoretical framework, termed the corporate stage, that synthesizes Erving Goffman's dramaturgical model of front-stage and back-stage performance, Michel Foucault's account of panoptic discipline, and Shoshana Zuboff's theory of surveillance capitalism to explain this reorganization. Where Goffman's back stage guaranteed performers a space in which to drop performance, and Foucault's panopticon relied on the mere possibility of observation to induce self-discipline, contemporary hybrid work collapses both guarantees: monitoring infrastructure operates continuously, disclosed rather than uncertain, converting workers' behavioral activity into extractable data. The chapter names and develops the resulting condition as the paradox of spectatorship: workers become simultaneously audience and performer, watching a datafied reflection of themselves while performing for an audience that is part algorithmic and part managerial, often without a stable way of knowing which. Drawing on recent work on worker-side sousveillance, the chapter also considers the limits of this paradox, arguing that opacity functions as a scarce and deliberately managed resource rather than a residual privacy the worker simply lacks. The chapter contributes to workplace and organizational studies a vocabulary for distinguishing platform-based visibility from its classical disciplinary predecessors, with implications for how monitoring systems might be designed without collapsing observation into total exposure.

\vspace{0.5em}
\noindent\textbf{Keywords:} algorithmic surveillance, hybrid work, dramaturgical theory, panopticism, surveillance capitalism, spectatorship, workplace visibility
\end{abstract}

\begin{callout}{Civil-Rights Framing.}
Although this chapter is rooted in critical theory rather than doctrine, its arguments bear directly on civil rights. Three links structure that fit: (1) differential impact --- surveillance's costs fall unevenly on workers lacking spatial privilege, along lines of class, gender, race, and disability; (2) equality and recourse --- gig-platform rating systems leave workers least able to contest an account of their conduct they did not author; and (3) opacity as a civil-rights interest --- a worker's capacity to remain occasionally unreadable deserves protection on par with legitimate managerial interests. The chapter's closing call for comparative regulatory analysis is the natural bridge from this theory to law.
\end{callout}

\section{Introduction}

A worker logs into a video call from a spare bedroom, angles the camera to exclude an unmade bed, and opens a second window running software that logs keystrokes, application switches, and idle time. Somewhere on a dashboard she will likely never see in full, these signals resolve into a productivity score. She is, in this moment, watched by a colleague on the call, watched by software she cannot see operating, and watching herself --- adjusting posture, tone, and visible activity in anticipation of both audiences at once. This scene, unremarkable in its ordinariness across contemporary hybrid and remote workplaces, is the empirical starting point for this chapter's argument: that the shift to hybrid and algorithmically monitored work has not simply added new tools to an old practice of supervision, but has reorganized the basic conditions under which workers are seen, and see themselves, at work.

Workplace visibility has never been neutral. What has changed is its architecture. Where earlier regimes of supervision relied on a supervisor's line of sight, a factory floor's layout, or a manager's sporadic presence, contemporary hybrid work distributes visibility across software that operates continuously, asymmetrically, and largely outside the worker's control. The question this chapter takes up is not whether workers are watched --- they always have been --- but how the specific combination of algorithmic monitoring and hybrid work reorganizes the relationship between watching and being watched, producing a condition this chapter calls the paradox of spectatorship.

\subsection{The Corporate Stage as Analytic Frame}

This chapter borrows, and deliberately strains, Erving Goffman's dramaturgical vocabulary. Goffman's front stage and back stage described a world in which performances were bounded in space and time: a worker performed professionalism on the office floor and relaxed once out of view. Hybrid work collapses this architecture. The home becomes a front stage that a webcam or activity monitor can activate at any moment, and the back stage --- the space in which a worker could once be off duty, off script, unobserved --- shrinks correspondingly or disappears. The corporate stage, as this chapter develops the term, names a performance space that is no longer bounded by walls or shifts but is instead switched on and off by software, with the worker frequently unable to know which condition currently obtains. This is not, strictly, a condition Goffman's model cannot accommodate --- he was well aware that back regions could be breached, that boundary control could fail, and that a performer could be caught between fronts. But Goffman treated such breaches as contingent hazards of performance, exceptions requiring explanation. The corporate stage names the case in which that exception becomes the rule: a performance space in which boundary failure is not an accident the performer risks but the default condition software imposes.

\subsection{Contributions / Chapter Roadmap}

The chapter makes the following moves:
\begin{itemize}[leftmargin=1.5em]
\item Synthesizes Goffman, Foucault, and Zuboff into a single framework for reading hybrid workplace visibility.
\item Names and develops the paradox of spectatorship as a distinct condition of contemporary algorithmic management.
\item Considers what this reorganization of visibility implies for worker autonomy and wellbeing, and for how organizations might design monitoring practices that do not collapse the distinction between observation and total exposure.
\end{itemize}

\section{Theoretical Framework}

\subsection{Goffman: The Presentation of Self as Front-Stage Labor}

Goffman's dramaturgical model treats social life as a series of performances staged for particular audiences, sustained through what he called impression management: the continuous, often unconscious work of controlling how one appears to others (Goffman). Central to this model is the distinction between front stage, where a performance is delivered, and back stage, where a performer can drop the performance, rehearse, or simply rest. Crucially, Goffman located the back stage in physical space --- a break room, a car, a home --- precisely because performances required somewhere to end.

Goffman's typology of roles that complicate the performer/audience division is also useful here, particularly the figure of the informer: someone granted access to the back region who nonetheless communicates what occurs there to an outside party the performer cannot fully monitor in turn (Goffman). Monitoring software occupies structurally the same position --- present in what should be back-region space, but oriented toward an audience the worker cannot see or negotiate with --- with the crucial difference that the informer is no longer a person capable of discretion, loyalty, or being persuaded, but a fixed technical function.

Goffman's dramaturgical vocabulary continues to be extended well beyond its original workplace and institutional settings; recent scholarship, for instance, has applied his account of stigma management and identity negotiation to autobiographical illness narrative (Almansoori and Amjad), underscoring the continued reach of his framework across otherwise unrelated domains of self-presentation.

Hybrid work erodes this spatial guarantee. When the home itself becomes a monitored workspace, the back stage does not merely shrink; it becomes conditional, switched on and off by whether a camera is active or a monitoring application is running. Workers report adapting home environments --- posture, lighting, even the visible contents of a room --- not for a single scheduled performance but for a front stage that could, in principle, be activated at any moment. Impression management, in this setting, becomes a standing condition of labor rather than an episodic task bounded by shift or meeting.

\subsection{Foucault: Panopticism and the Internalization of the Gaze}

Foucault's account of the panopticon turns on a specific mechanism: disciplinary power that operates not through constant observation but through the uncertainty of being observed, such that the watched party internalizes the gaze and disciplines their own conduct regardless of whether anyone is actually watching at a given moment (Foucault). This uncertainty principle is what makes the panopticon efficient --- surveillance does not need to be continuous to be effective, only possible. Algorithmic workplace monitoring both intensifies and departs from this model. It intensifies it insofar as software can, in principle, observe continuously, removing the gaps in which a watched party could once assume privacy by default. But it also departs from Foucault's model in a specific way: where the panopticon relies on uncertainty to produce self-discipline, many monitoring systems make their presence explicit --- a recording indicator, a dashboard the worker can also see --- converting internalized discipline into an openly acknowledged, sometimes actively managed, condition. The worker under algorithmic monitoring is often not guessing whether they are watched, but calibrating behavior against a known, quantified record of having been watched.\footnote{Concerns of this kind are not confined to the historical or workplace contexts discussed here; the same uncertainty-based logic of surveillance surfaces in other institutional settings, including recent accounts of political organizing on university campuses, where participants weighed decisions about physical space partly around fear of being surveilled (Iyer and D'Souza).}

\subsection{Zuboff: Surveillance Capitalism and Behavioral Data Extraction}

Zuboff's theory of surveillance capitalism describes an economic logic in which human experience is unilaterally claimed as raw material, translated into behavioral data, and rendered into predictive products sold in markets the person generating the data never sees or consents to in any meaningful sense (Zuboff). Zuboff developed this account primarily through consumer platforms --- search, social media, smart devices --- but the underlying logic extends with little modification into the workplace. Keystroke logs, activity timestamps, meeting attendance patterns, and productivity scores constitute a comparable behavioral surplus, extracted not from consumers but from labor itself, and often justified internally as a byproduct of legitimate productivity or security tools rather than named as data extraction. The asymmetry Zuboff identifies as central to surveillance capitalism --- one party accumulates knowledge and predictive power, the other supplies the raw behavioral material largely without recourse --- maps closely onto the relationship between an employee and the monitoring infrastructure of a hybrid workplace.

\subsection{Toward a Synthesis: Hybrid Visibility}

Read together, these three accounts describe three distinct but compounding operations: Goffman locates the site of performance and shows how that site has lost its guaranteed off-switch; Foucault explains the disciplinary mechanism by which observation, actual or possible, shapes conduct; and Zuboff supplies the economic logic that gives monitoring its institutional purpose --- behavioral data is not merely observed but captured, aggregated, and rendered actionable. Hybrid visibility, as this chapter uses the term, names the condition produced when these three operations occur together: a performance space without a reliable back stage (Goffman), sustained by infrastructure that watches continuously rather than merely possibly (an intensification of Foucault), for the purpose of producing extractable behavioral data (Zuboff). This convergence is not confined to workplace software: comparable dynamics appear in adjacent domains where dashboards render a person's activity legible to both themselves and an institution, as in learning analytics dashboards designed to make distance learners' study behavior visible to instructors and, reflexively, to the learners themselves (Rets et al.). The design problem such dashboards pose --- what to surface, to whom, and with what effect on the surveilled party's self-perception --- is structurally the same problem posed by workplace productivity dashboards, suggesting that hybrid visibility is not unique to employment but is one instance of a broader pattern in dashboard-mediated institutional relationships. None of the three theorists anticipated this specific convergence, which is precisely what makes their combination useful: each supplies a piece the others lack.

\section{Hybrid Visibility in the Modern Workplace}

\subsection{From Panopticon to Platform: What Changes}

Classical disciplinary visibility, in Foucault's account, was architectural: it depended on a specific arrangement of walls, sightlines, and thresholds, such that a guard tower, a factory floor's open layout, or a manager's office positioned to overlook a workroom did the disciplinary work of shaping conduct. Its efficacy rested on uncertainty --- the watched party could not confirm, at any given moment, whether the gaze was actually present --- and on impermanence: even when observation occurred, it typically left no durable record. A supervisor's glance across a factory floor disciplined the moment it occurred but did not, in itself, become data that could be replayed, aggregated, or compared against a worker's conduct the following week.

Platform-based visibility departs from this architecture along three axes. First, it is continuous rather than intermittent: monitoring software does not need to be physically present to observe, and it does not tire, look away, or leave the room. Second, it is quantified rather than impressionistic: where a supervisor's judgment of a worker's effort was qualitative and contestable, an activity score or keystroke count presents itself as an objective measurement, even when the underlying proxy --- keystrokes, mouse movement, application-switching frequency --- bears an uncertain relationship to actual productivity or quality of work. Third, and most consequentially, it is retentive: platform-based monitoring does not merely observe in the moment but stores what it observes, producing a durable record that can be aggregated over days or months, compared across workers, and retrieved for review long after the conduct it describes has ended. This retentiveness is what Foucault's model, built around momentary uncertainty, could not anticipate: contemporary workplace monitoring converts the disciplinary gaze from something a worker internalizes and eventually stops attending to into something that persists as an external, queryable record independent of the worker's own memory or account of events.

The effect of this shift is to relocate the site of discipline. Where architectural surveillance disciplined conduct in the moment through the possibility of being seen, platform-based visibility disciplines conduct retrospectively, through the certainty of having been recorded. A worker under the panopticon asks, in effect, ``am I being watched right now?'' A worker under hybrid visibility asks a different question: ``what will this data say about me later, to an audience I cannot fully anticipate?'' That question --- oriented toward a future, cumulative judgment rather than a present, uncertain one --- is the specific innovation that platform-based visibility introduces, and it is what makes the corporate stage, as this chapter uses the term, a genuinely distinct condition rather than simply a more efficient panopticon.

\subsection{The Home as Front Stage}

The domestic space was, for Goffman, an almost paradigmatic back stage: the place to which a performer retreated to remove a costume, drop an accent, or simply stop managing an impression. Hybrid work reassigns this space a second, competing function. The same room in which a worker eats, sleeps, or cares for a child becomes, for the duration of a video call or an active monitoring session, a professional front stage --- one that must be composed, lit, and framed for an audience, and that can be activated with little warning by a calendar invitation or an automatically triggered webcam check-in. Workers describe the resulting adjustments in strikingly theatrical terms even when they do not reach for Goffman's vocabulary themselves: repositioning a camera to exclude an unmade bed or a child's toys, keeping a jacket on a chair within reach for an unscheduled call, or maintaining a tidy, camera-ready corner of an otherwise lived-in room.

This spatial collapse is compounded by a second, distinct one. Goffman's model assumes performers can maintain audience segregation: sustaining different fronts for different audiences by ensuring those audiences never occupy the same space at the same time --- the version of oneself performed for coworkers need not be reconciled with the version performed for family, so long as the two never watch simultaneously. A video call breaches this segregation directly. A child wandering into frame, a partner's voice audible in the background, or a housemate visible through a doorway does not merely risk an awkward moment; it merges audiences Goffman's model treats as requiring separation, forcing a single front to serve simultaneously as both. Where the spatial argument above concerns whether a back stage exists at all, audience segregation names a related but separate failure: even when a back region nominally remains, hybrid work removes the wall that once kept its audience from mixing with the front stage's. These are not incidental accommodations but a form of standing stagecraft, performed in anticipation of a front stage that could, in principle, switch on at any moment.

What distinguishes this arrangement from Goffman's original model is that the switch between front and back stage is no longer under the performer's control. An office worker leaving a meeting could reliably step into a hallway, a break room, or a car and know that the performance had ended. A hybrid worker's transition out of performance is instead conditional on software: whether a monitoring application is currently active, whether a webcam indicator light is lit, whether an activity tracker is still logging keystrokes after a call has formally ended. The back stage does not disappear so much as become intermittent and administratively determined, available only in the gaps the monitoring infrastructure happens to leave open.

This reorganization is not evenly distributed.

\begin{callout}{Civil-Rights Mapping.}
This claim is the chapter's strongest point of contact with civil-rights scholarship: the collapse of front/back stage is not evenly distributed, falling heaviest on workers without spatial privilege --- mediating inequalities along lines of class, gender, race, and disability. That differential exposure is precisely the kind of disparate-impact and non-discrimination question civil-rights doctrine exists to address.
\end{callout}

Workers with a dedicated home office can install a durable boundary --- a closed door, a room used for no other purpose --- that approximates Goffman's spatial back stage even under monitoring. Workers without that spatial privilege, whether because of housing size, shared occupancy, or caregiving responsibilities that make a room's non-professional use unavoidable, absorb the cost of the corporate stage's expansion directly into domestic space that was never designed to sustain a professional performance at all. The collapse of front and back stage under hybrid visibility is, in this sense, not only a psychological or disciplinary phenomenon but a material one, mediated by the same inequalities in housing and domestic labor that structure hybrid work more broadly.

\subsection{Metrics, Scores, and the Datafied Self}

If the home becomes a front stage, the dashboard becomes its mirror. Productivity software translates a worker's activity --- keystrokes, application windows, idle time, meeting attendance --- into a composite score, an active time percentage, or a ranked comparison against peers. This is Zuboff's behavioral-surplus logic operating at the scale of an individual employee: raw conduct is rendered into a data product, and that product circulates to an audience --- a manager, an HR system, sometimes the worker themselves --- that did not observe the underlying conduct directly and knows it only through its datafied form.

This register of visibility is distinct in kind from direct human observation, not merely more efficient than it. A supervisor watching a worker type could recognize context: a long pause might register as thinking, reading, or a difficult problem rather than as idleness. A keystroke logger cannot make this distinction; it records absence of input as absence of activity, regardless of cause, and renders that absence as a quantified deficit on a dashboard. The worker is thus not simply watched more, but watched differently, through a medium that strips out precisely the contextual judgment a human observer could supply. Workers respond, predictably, by managing the metric rather than only the underlying conduct it purports to measure: moving a mouse periodically to avoid being flagged as idle, keeping an application window open regardless of whether it is in active use, or timing visible activity to coincide with periods the worker believes are more heavily reviewed. This is not evidence of a workforce being disciplined into diligence; it is evidence of a workforce learning to perform for a legible proxy rather than for the underlying goal that proxy was meant to approximate --- a familiar problem in measurement generally, but one that acquires particular force when the proxy is also, per Section 4, something the worker is made to watch.

Goffman's vocabulary supplies a name for the effort this requires: dramaturgical discipline, the sustained self-control needed to carry out a performance without a revealing slip --- maintaining presence of mind, suppressing an inappropriate reaction, keeping minor errors from disrupting the impression being conveyed (Goffman). Goffman located this discipline in the performer's relation to a human audience capable of noticing a lapse. What has been described above is the same discipline redirected at an audience that cannot be reasoned with or read: a worker moving a mouse to avoid an idle flag is not managing an impression for someone but performing dramaturgical discipline at a proxy incapable of registering intent, context, or explanation --- only presence or absence of signal. This is a narrower, harder task than the one Goffman described, since a human audience can be reassured, and a keystroke logger cannot.

The datafied self this produces is, crucially, not identical to the worker's own sense of their conduct. A worker who spent an hour in careful, screen-away thought about a problem and a worker who spent that hour genuinely idle can be rendered indistinguishable by the same dashboard, collapsing a meaningful difference in labor into an identical data point. The metric becomes, in effect, a second, competing account of the worker's day --- one that the worker did not author, cannot always contest, and must nonetheless treat as authoritative, since it is this account, not the worker's own recollection, that circulates to managerial and algorithmic audiences. This is the material infrastructure that makes the paradox of spectatorship (Section 4) possible: the worker becomes an audience to their own conduct only because that conduct has already been converted into a form --- the score, the dashboard, the datafied self --- that can be watched at all.

\section{The Paradox of Spectatorship}

\subsection{Audience and Performer at Once}

Before this condition can be named, it must first be felt as an absence. Workers describing algorithmic monitoring rarely reach immediately for a term as precise as surveillance or as theoretically loaded as panopticon. What they report first is closer to an unlocatable discomfort: not fear of a supervisor's gaze, which is at least addressed to someone, but a flatter and harder-to-place unease at being rendered into data by a system with no face and no obvious audience on the other end. This is worth pausing on, because it names precisely the kind of theoretical problem that Sianne Ngai, writing on Fredric Jameson's method of argumentation, has recently described as characteristic of concepts still under construction: a moment in which the existing name for a cultural phenomenon --- the concept we automatically reach for first --- suddenly seems insufficient, leaving in its place the absence of a concept ready to fill its place (Ngai). Ngai's account of Jameson's interpretive procedure --- feeling, to noncognition, to a felt experience of cognition --- offers a useful diagnostic for the paradox this chapter is naming. Goffman's audience presupposes a person; Foucault's guard in the tower presupposes, at minimum, a tower and an inmate who can imagine someone in it. What algorithmic workplace monitoring produces is visibility without a stable witness: the sensation of being watched, detached from any confident answer to the question of by whom, or even by what. The discomfort this produces is not simply an intensification of older surveillance anxieties; it is, in Ngai's sense, a nonfeeling --- the felt absence of the expected feeling of being seen by someone who could, in principle, see you back. Reader-response theory offers a useful point of contrast here. In that tradition, a text's meaning is understood as something completed by an active, if abstract, reader --- Iser's ``implied reader,'' invoked in recent reader-response scholarship on autobiographical illness narrative to describe how a memoir's gaps are filled in by the reader's own experience (Vailankanni). That model, however abstracted the reader, still presupposes an interpreting consciousness on the other end of the exchange. The algorithmic audience this chapter describes offers no such consciousness to complete anything: it registers signal, not meaning, which is precisely what leaves the worker's sense of being watched unanchored to any imaginable watcher.

This is the gap that the paradox of spectatorship is meant to fill. The worker under hybrid visibility occupies two positions that Goffman's model treats as sequential and Foucault's model treats as external to the watched party, but which algorithmic monitoring collapses into simultaneity. First, the worker is a performer: adjusting camera angle, tone, and visible activity for an audience that is part managerial and part infrastructural, whose composition at any given moment the worker frequently cannot verify. Second, and at the same time, the worker is that performance's first audience: consulting the same productivity dashboard, activity score, or active time metric that management consults, watching a datafied reflection of their own conduct assembled by a system they did not design and cannot fully audit. These are not two workers taking turns; they are one worker holding both positions at once, in a loop where self-monitoring in anticipation of the algorithmic audience becomes indistinguishable from performing for it. Goffman's actor could step back stage and stop performing; Foucault's prisoner, uncertain whether the tower was occupied, could at least locate the direction from which discipline arrived. The worker under hybrid visibility has neither exit: no reliable back stage, per Section 3, and no single, locatable direction in which the gaze originates, since the gaze is distributed across software, aggregated metrics, and the worker's own anticipatory self-surveillance. The paradox of spectatorship names this specific, previously undertheorized convergence: not surveillance simply intensified, but spectatorship without a stable spectator, performance without a bounded stage, and self-knowledge that arrives, if it arrives at all, mediated through the same infrastructure that produces the exposure it is meant to help the worker manage.

\subsection{Illustrative Cases}

The paradox of spectatorship is easiest to see where it is most concentrated. Three settings illustrate the framework developed above at different points along the spectrum from human to algorithmic audience.

\textbf{Remote productivity software.} Activity-tracking applications, now standard in many hybrid and remote workplaces, generate the clearest instance of the datafied self described in Section 3. These tools log input activity and translate it into an active time figure visible on a shared dashboard, often accessible to the worker in the same form it reaches their manager. The worker is thereby positioned to monitor their own score in anticipation of a managerial review that may or may not occur on any given day --- a textbook instance of the simultaneity described above, in which watching one's own metric and performing for the audience that will eventually see it become a single, continuous activity rather than two sequential ones.

\textbf{Call center monitoring.} Call centers represent an older, partially algorithmic register of workplace visibility, and are useful precisely because they show the corporate stage in an intermediate form. Metrics such as average handle time and, more recently, automated sentiment or tone analysis applied to recorded calls combine real-time human supervision (a manager listening in) with algorithmic scoring layered on top of it. Workers in this setting are visible to two audiences of different character at once --- a human supervisor whose attention is intermittent, and a scoring system whose attention is total --- reproducing, in a setting that predates most contemporary hybrid-work software, the same uncertainty about which audience is currently active that this chapter identifies as central to the paradox of spectatorship.

\textbf{Gig-platform rating systems.} Ride-share and delivery platforms extend algorithmic visibility furthest from any human supervisor. A driver's or courier's conduct is rendered legible almost entirely through customer ratings, GPS-tracked route efficiency, and platform-side acceptance-rate metrics, aggregated into a score that determines continued access to work. There is, in this setting, no back stage in even the conditional sense described in Section 3: the vehicle or delivery route is simultaneously workplace and monitored surface for the duration of a shift, and the managerial audience is, for practical purposes, entirely algorithmic. This is also the setting in which the sousveillance tactics discussed below --- informal worker networks such as Turkopticon-style reputation tools --- have been most fully documented, since gig workers facing an almost purely algorithmic audience have the least recourse to the human, negotiable channels of contestation available in more conventionally supervised workplaces.

Read together, these three cases suggest that the paradox of spectatorship is not a binary condition a workplace either has or lacks, but a spectrum defined by the proportion of algorithmic to human audience and by how much conditional back stage, if any, remains available to the worker.

\subsection{Limits of the Paradox: Resistance and Opacity}

The paradox of spectatorship is not a closed system, and workers are not simply its passive subjects. A growing body of research on worker-side monitoring --- termed sousveillance, the inversion of surveillance in which those observed turn monitoring back on those in power --- shows that workers actively develop tactics to reclaim some of the asymmetry hybrid visibility produces. In a study of gig workers, De Los Santos et al. found that workers routinely research and document information about their employers and platforms: tracking payment discrepancies, cross-referencing requester reputations through informal networks and tools like Turkopticon, and identifying gaps in available information as a signal in themselves (De Los Santos et al.). These tactics echo Mathiesen's reading of Foucault, in which sousveillance is theorized as a structural inversion of the panopticon's gaze --- the watched turning documentation back on the watcher rather than merely enduring it.

Crucially, this resistance is neither total nor uncomplicated. Workers in that study weighed sousveillance against real costs: the time and effort required, the risk of damaging a working relationship, contractual and platform restrictions on collecting data about employers, and the emotional toll of what such information sometimes reveals. Several participants described choosing not to look further once they had enough to protect themselves, precisely because further looking carried relational or psychological cost. This suggests that opacity --- a worker's capacity to remain partially unreadable to a monitoring system, or to choose not to fully exploit the visibility available to them --- functions as its own limited resource, exercised deliberately rather than simply imposed from outside. The paradox of spectatorship, in other words, does not eliminate worker agency; it reshapes the terms on which that agency is exercised, requiring workers to manage not only their own visibility but their own capacity to see.

The concept of sousveillance itself predates its application to gig work. Steve Mann coined the term to describe an inversion of surveillance --- watching ``from below'' rather than from above --- theorizing it as a kind of inverse panopticon in which recording and monitoring capacity is redistributed to those ordinarily subject to observation rather than reserved to those in authority over them (Mann, Nolan, and Wellman). Read against that genealogy, the tactics De Los Santos et al. document among gig workers are a specific, labor-inflected instance of a broader pattern: sousveillance as a structural response to asymmetric visibility wherever it occurs, of which the algorithmically monitored workplace is one site rather than the origin.

\section{Discussion}

The preceding sections have argued that hybrid, algorithmically monitored work does not merely intensify older forms of workplace supervision but reorganizes their basic conditions, producing a paradox of spectatorship in which the worker is simultaneously performer and audience for a datafied account of their own conduct. This section considers what that reorganization implies for organizational design and management practice, and situates the chapter's contribution within this volume's broader concern with the evolving challenges and opportunities of the contemporary workplace.

The most direct implication is that opacity, as Section 4 argued, should be treated by organizations as a resource to be deliberately allocated rather than a residual gap to be closed through ever more complete monitoring. Systems designed on the assumption that more visibility is unambiguously better --- more granular activity logs, more continuous scoring, more dashboards available to more parties --- risk deepening the paradox of spectatorship rather than managing it, by further collapsing the conditional back stage identified in Section 3 and increasing the burden of self-surveillance this chapter has described. Organizations seeking to design monitoring systems that do not collapse observation into total exposure might instead ask, of any new monitoring capability, whether it restores some locatable back stage to the worker --- a defined window in which monitoring is reliably inactive, a boundary between data collected for coordination and data collected for evaluation, a clear and stable answer to the question of who or what is watching at a given moment --- rather than treating comprehensiveness of visibility as an end in itself.

The chapter's framework also has implications for worker wellbeing that extend beyond conventional accounts of surveillance-related stress. Because the paradox of spectatorship positions the worker as their own audience, the psychological cost it describes is not reducible to the discomfort of being watched by another person; it includes the additional labor of continuously interpreting a datafied reflection of oneself that may not correspond to one's own sense of the work performed, as Section 3 argued in discussing the gap between a worker's account of their day and its algorithmic record. Management practices that treat activity scores as self-evidently accurate, without acknowledging this gap, risk compounding rather than resolving that cost, particularly for workers --- as Section 3 also noted --- whose domestic circumstances make a stable front-stage/back-stage boundary hardest to maintain in the first place.

Positioned within this volume's broader concern with the evolving landscape of workplace challenges and opportunities, the paradox of spectatorship contributes a vocabulary for a specific kind of challenge: one that arises not from any single monitoring technology in isolation, but from the compounding effect of continuous, quantified, retained observation layered onto work environments --- the home chief among them --- that were never designed to sustain it. The corresponding opportunity, this chapter suggests, lies less in better monitoring technology than in organizational designs that treat the worker's capacity to remain occasionally unreadable as a legitimate design constraint, on par with the legitimate managerial interest in visibility that motivates monitoring in the first place.

\section{Conclusion}

This chapter has argued that the shift to hybrid, algorithmically monitored work reorganizes the conditions under which workers are seen and see themselves, producing a condition it has named the paradox of spectatorship: the worker as simultaneously performer and first audience for a datafied account of their own conduct, assembled by infrastructure they did not design and cannot fully audit. Building the corporate stage as an analytic frame --- a synthesis of Goffman's dramaturgical model, Foucault's account of disciplinary visibility, and Zuboff's theory of behavioral-data extraction --- has made it possible to identify what is genuinely new in this arrangement rather than treating it as a straightforward intensification of familiar workplace supervision. Where Goffman's back stage guaranteed an off-switch, Foucault's panopticon relied on uncertainty, and Zuboff's surveillance capitalism operated primarily on consumers rather than employees, hybrid visibility combines a conditional, software-mediated back stage with continuous, disclosed observation and the systematic conversion of labor into extractable data --- a convergence none of the three theorists, working separately, had occasion to anticipate. At the same time, as Section 4 argued through the case of worker-side sousveillance, this condition does not eliminate worker agency so much as redistribute it, making opacity itself --- the capacity to remain partially unreadable, or to choose not to look too closely at what a monitoring system reveals --- a scarce resource that workers manage deliberately rather than simply lack.

The framework developed here opens several directions for further research that this chapter's scope has not allowed it to pursue.

\begin{callout}{Bridge to Law.}
This comparative-regulatory direction is the explicit pivot from critical theory to law: how differing jurisdictional protections for employee data and electronic surveillance shape the intensity of the paradox, and whether disclosure regimes restore a meaningful ``back stage'' or merely formalize its absence.
\end{callout}

Comparative regulatory work could examine how differing legal approaches to workplace monitoring and algorithmic management --- across jurisdictions with markedly different baseline protections for employee data and electronic surveillance --- shape the intensity of the paradox of spectatorship in practice, and whether disclosure requirements of the kind already discussed in Section 2 are sufficient to restore a meaningful back stage or merely formalize its absence. Existing statutes offer an instructive contrast in regulatory philosophy. In the United States, a small number of states --- New York, Connecticut, and Delaware among them --- require only that employers disclose electronic monitoring in advance, without restricting what conclusions employers may draw from the resulting data or how those conclusions may be used (New York State; State of Connecticut). Connecticut's statute has recently been amended to require that disclosure specify not merely that monitoring occurs but where on the employer's premises it occurs, an incremental move toward exactly the kind of locatable back stage this chapter's Discussion section calls for, though the amendment still stops short of limiting what may be done with the data once collected (State of Connecticut). The European Union's General Data Protection Regulation takes a structurally different approach: Article 22 restricts not the act of monitoring but a downstream consequence of it, establishing a qualified right not to be subject to decisions based solely on automated processing or profiling that produce legal or similarly significant effects, together with a right to request human intervention (European Union). Where the American disclosure model addresses the worker's right to know that the corporate stage is active, the European model addresses the worker's right to contest what the stage's audience does with what it records --- a distinction that maps closely onto this chapter's argument that disclosure alone, of the kind discussed in Section 2.2, converts uncertainty into an openly acknowledged condition without thereby limiting what the resulting record may be used to justify. Further empirical attention to worker-side sousveillance tools, extending the study discussed in Section 4 beyond gig work into conventionally salaried hybrid employment, would help clarify whether the tactics available to gig workers generalize to settings with more stable employment relationships and correspondingly different costs and protections. Finally, the design implications raised in Section 5 --- monitoring systems built around a deliberately preserved back stage rather than maximal visibility --- warrant testing against actual organizational practice, to determine whether opacity-by-design is a workable management principle or remains, for now, primarily a corrective vocabulary for diagnosing what current systems lack.

\newpage
\begin{thebibliography}{99}
\setlength{\itemsep}{0.6em}

\bibitem{almansoori} Almansoori, S. T. K., and F. A. Amjad. ``Negotiating Stigma and Reconstructing Identity in Matt Haig's \emph{Reasons to Stay Alive}: A Goffmanian Reading''. \emph{Journal of Daoist Studies} 19-B (2026): 1414--1427.

\bibitem{delossantos} De Los Santos, M., et al. ``Designing Sousveillance Tools for Gig Workers''. Northeastern University, IBM Research, and Universidad Nacional Aut\'onoma de M\'exico, arXiv preprint arXiv:2403.09986 (2024).

\bibitem{eu} European Union. Regulation (EU) 2016/679 (General Data Protection Regulation), Article 22: Automated Individual Decision-Making, Including Profiling, 2016.

\bibitem{foucault} Foucault, Michel. \emph{Discipline and Punish: The Birth of the Prison}. English translation by Alan Sheridan, Vintage Books, 1977. \'Editions Gallimard, 1975.

\bibitem{goffman} Goffman, Erving. \emph{The Presentation of Self in Everyday Life}. Garden City, NY: Doubleday, 1959.

\bibitem{iyer} Iyer, Usha, and Sophie D'Souza. ``Decolonial Pedagogy and Praxis in a Time of Live-Streamed Genocide: Dialogic Notes from a Global North Classroom''. In \emph{Decolonial Entanglements: Praxis, Pedagogy, and Social Theory}, ed. by Jairo I. F\'unez-Flores, Ana Carolina D\'iaz Beltr\'an, and Nathalia E. Jaramillo, 78--94. New York: Routledge, 2026.

\bibitem{mann} Mann, Steve, Jason Nolan, and Barry Wellman. ``Sousveillance: Inverse Surveillance in Multimedia Imaging''. In \emph{Proceedings of the 12th Annual ACM International Conference on Multimedia}, 620--627. New York: ACM, 2004. \url{https://doi.org/10.1145/1027527.1027673}.

\bibitem{nystate} New York State. Civil Rights Law \S 52-c: Electronic Monitoring of Employees. Effective May 7, 2022, 2022.

\bibitem{ngai} Ngai, Sianne. ``Thresholds of the Concept''. In \emph{The Future of Totality: Fredric Jameson and the Prospects of Critical Theory}, ed. by Nicholas Brown et al., 17--27. Durham, NC: Duke University Press, 2026.

\bibitem{rets} Rets, Ia, et al. ``Exploring Critical Factors of the Perceived Usefulness of a Learning Analytics Dashboard for Distance University Students''. \emph{International Journal of Educational Technology in Higher Education} 18, no. 1 (2021): Article 46.

\bibitem{ctstate} State of Connecticut. Connecticut General Statutes \S 31-48d: Electronic Monitoring of Employees. As amended by Public Act No. 26-73 (2026), effective October 1, 2026, 1998.

\bibitem{vailankanni} Vailankanni, M. M. C. ``Reconstructing Mental Health Narratives: A Reader Response Approach to \emph{Reasons to Stay Alive} by Matt Haig''. \emph{International Journal of Literature, Linguistics, and Humanities} 1, no. 2 (2025): 20--23. \url{https://doi.org/10.64137/31078729/IJLLH-V1I2P105}.

\bibitem{zuboff} Zuboff, Shoshana. \emph{The Age of Surveillance Capitalism: The Fight for a Human Future at the New Frontier of Power}. New York: PublicAffairs, 2019.

\end{thebibliography}

\newpage
\appendix
\section*{Appendix: Empire Under the Night Sky --- HoST Dossier Citations}
\addcontentsline{toc}{section}{Appendix: Empire Under the Night Sky --- HoST Dossier Citations}

\textit{The following bibliography is unrelated to this chapter's argument. It is appended here at the author's request and reproduces the citation set for the thematic dossier ``Empire Under the Night Sky: Recording Astral-Cosmography in Qing Dynasty China, 17th--19th Centuries,'' \emph{HoST -- Journal of History of Science and Technology} 18, no.\ 1 (June 2024).}

\subsection*{Dossier Articles}

\begin{thebibliography}{99}
\setlength{\itemsep}{0.6em}

\bibitem{brown2024introduction} Brown, Tristan G. ``Introduction: Empire under the Night Sky: Recording Astral-Cosmography in Qing Dynasty China, 17th--19th Centuries.'' \emph{HoST -- Journal of History of Science and Technology} 18, no. 1 (June 2024): 1--5. \url{https://doi.org/10.2478/host-2024-0001}.

\bibitem{chen2024fenye} Chen, Shih-Pei. ``Fenye by the Numbers: A Quantitative Analysis of Astrological Contents in Chinese Local Gazetteers.'' \emph{HoST -- Journal of History of Science and Technology} 18, no. 1 (June 2024): 6--30. \url{https://doi.org/10.2478/host-2024-0002}.

\bibitem{wu2024encounter} Wu, Huiyi. ``An Encounter of Incommensurables: European Cosmological Knowledge in the Fenye Chapters of Chinese Local Gazetteers (1660--1820).'' \emph{HoST -- Journal of History of Science and Technology} 18, no. 1 (June 2024): 31--60. \url{https://doi.org/10.2478/host-2024-0003}.

\bibitem{brown2024fengshui} Brown, Tristan G. ``From Fenye to Fengshui: Applying Correlative Cosmography in Late Imperial China.'' \emph{HoST -- Journal of History of Science and Technology} 18, no. 1 (June 2024): 61--85. \url{https://doi.org/10.2478/host-2024-0004}.

\bibitem{zhang2024latitude} Zhang, Jiajing. ``From Fenye to Latitude and Longitude: The Adjustment and Renewal of the Positioning System in Chinese Local Gazetteers during Qing Dynasty.'' \emph{HoST -- Journal of History of Science and Technology} 18, no. 1 (June 2024): 86--114. \url{https://doi.org/10.2478/host-2024-0005}.

\end{thebibliography}

\subsection*{Works Cited in Brown's Introduction (Footnotes)}

\begin{thebibliography}{99}
\setlength{\itemsep}{0.6em}

\bibitem{weil2022chinese} Weil, Dror. ``Chinese-Muslims as Agents of Astral Knowledge in Late Imperial China.'' In \emph{Overlapping Cosmologies in Asia: Transcultural and Interdisciplinary Approaches}, edited by Bill M. Mak and Eric Huntington, 116--138. Leiden: Brill, 2022.

\bibitem{pankenier2005characteristics} Pankenier, David W. ``Characteristics of Field Allocation (\emph{fenye}) Astrology in Early China.'' In \emph{Current Studies in Archaeoastronomy: Conversations Across Space and Time}, edited by John Fountain and Rolf M. Sinclair, 499--513. Durham: Carolina Academic Press, 2005.

\bibitem{elman2005ownterms} Elman, Benjamin. \emph{On Their Own Terms: Science in China, 1550--1900}. Cambridge, MA: Harvard University Press, 2005.

\bibitem{ge2018intellectual} Ge, Zhaoguang. \emph{An Intellectual History of China, Volume Two}. Translated by Josephine Chiu-Duke and Michael S. Duke. Leiden: Brill, 2018.

\bibitem{qiu2020tiandi} Qiu, Jingjia. \emph{Tiandi zhi jian: Tianwen fenye de lishixue yanjiu} [Between Heaven and Earth: A Historical Study of Fenye Astrology]. Beijing: Zhonghua shuju, 2020.

\bibitem{tackett2017origins} Tackett, Nicolas. \emph{The Origins of the Chinese Nation: Song China and the Forging of an East Asian World Order}. Cambridge: Cambridge University Press, 2017.

\bibitem{chen2023logart} Chen, Shih-Pei, and Calvin Yeh. \emph{LoGaRT: Local Gazetteers Research Tools} (software). Berlin: Max Planck Institute for the History of Science, 2023. \url{https://logart.mpiwg-berlin.mpg.de/}.

\end{thebibliography}

\end{document}
